\begin{figure}[tb]
\centering
\footnotesize
\begin{minted}[
  xleftmargin=2ex, xrightmargin=1ex, frame=lines, fontsize=\footnotesize, linenos, 
  autogobble, numbersep=1ex, framesep=1mm, breaklines]
{java}
for (int i = 0; i < 10; i++) { // Loop 1
  sock.read(); // TP1
  if (condition1) {// TP2; MP1
    throw new IOException();
  }
  for (int j = 0; j < 10; j++) { // Loop 2
    if (foo()) { // TP3; MP2
      throw new IOException();
    }
  }
}

boolean foo() { // NP1
  fis = new FileInputStream(path); // TP 4
  if (fis.read() == -1) { // MP 3; TP 5
    return false;
  }
  return true;
}
\end{minted}
% \vspace{-3mm}
\caption{Pseudo-code demonstrating the injection and monitor points. ``TP'' means throw point and ``NP'' negation point. 
``MP'' means monitor point (used and explained in \Cref{sec:chain-forming:approx-check}).}
% \vspace{-1mm}
\label{fig:injection-point-exp}
\end{figure}
